\newcommand{\F}{\mathbb{F}}
\newcommand{\C}{\mathbb{C}}
\newcommand{\an}{\mathscr{A}} 
\newcommand{\Q}{\mathbb{Q}}
\renewcommand{\hom}{\operatorname{Hom}}
\newcommand{\ehom}{\operatorname{End}}
\colorlet{gris}{black!60}

\-\vspace{-0.3cm}
  {\hrule \vspace{-0.15cm}
  \centering \textbf{P1} \vspace{0.1cm}
  \hrule}

    \newcommand{\Z}{\mathbb{Z}}
\vspace{0.2cm}
1. (a) Defina subgrupo normal y diga cuales son los subgrupos normales de \(A_4\). 

\textcolor{gray}{
  \(N\leq G\) es \emph{subgrupo normal} si \(gNg^{-1} = N, \ \forall g \in G\). Recordando que \(|A_4| = \nicefrac{4!}{2} = 12\), por el \(\blacksquare\) \textbf{2.11. (\emph{Lagrange})} \(|N| \in \{1, 2, 3, 4, 6, 12\}\), analizaremos por casos:  
  \begin{itemize}[left = 0cm]
    \item \( 1 ;12 \ \mid \  \) Subgrupos normales triviales \(\{\id\} \) e \(A_4\).
    \item \(6\ \mid \ \) No tiene, subgrupo de índice \([G:N] = 2\) es normal pues \( G  = N \sqcup gN = N \sqcup Ng \ \leftrightharpoons \ gN = Ng\). Luego \((gN)^2 = N \ \leftrightharpoons \ g^2 \in N, \ \forall g \in G\), contiene todos los cuadrados de \(G\). En \(G = A_4\) los \(3\)-ciclos son cuadrados pues \(\alpha = (\alpha^2)^2 = \alpha^3 \alpha = \alpha\), hay \(8\) en total.\(\lightning\)  
    \item \(4 \ \mid \ \) \(N = \textbf{V}_4 = \{\id, (12) (34), (13) (24), (1 4)(23)\}\), no hay mas pues los restantes son \(3\)-ciclos de orden \(3 \nmid 4\).     
    \item \(3 \ \mid \ \) Generado por \(3\)-ciclo, sin embargo todos los \(3\)-ciclos son conjugados (por \(3\)-ciclos que están en \(A_4\)), 
    \[ (234)(123)(234)^{-1} = (134) \neq (123)^{-1},\]
    sigue que \(N \) no puede ser normal de orden \(3\). 
    \item \(2 \ \mid \ \) Los únicos elementos de orden \(2\) son producto de transposiciones disjuntas, todos en una misma clase:
    \[(123)\cdot(12)(34)\cdot (123)^{-1} = (23)(14).\] 
  \end{itemize}
    \vspace{-0.3cm}
  Por tanto, los subgrupos normales de \(A_4\) son \(\id, \ \textbf{V}_4\) e \(A_4\). 
}

(b) Encuentre todos los homomorfismos \(\varphi : A_4 \to \Z_{12}\). 

\textcolor{gray}{
  Sabemos que \(\ker \varphi \lhd A_4\), luego, del ítem (a) tenemos: 
  \begin{itemize}[left = 0cm]
    \item \(\ker \varphi = \id \ \mid \ \) \(\varphi^{-1}: \Z_{12}\cong A_4\) (no abeliano).\(\lightning\).
    \item \(\ker \varphi = \textbf{V}_4 \ \mid\  \) Por la propiedad universal del espacio cociente, siendo \(\pi : A_4 \to \nicefrac{A_4}{\textbf{V}_4} \cong \Z_3\) tenemos que \(\forall \varphi : A_4 \to \Z_{12}\)  homomorfismo \(\exists ! \widetilde{\varphi} : \nicefrac{A_4}{\textbf{V}_4} \to \Z_{12}\) homomorfismo \(: \varphi = \widetilde{\varphi} \circ \pi\) de los cuales hay apenas \(3\): 
    \[\widetilde{\varphi} : 1 \mapsto a \in \{1 , 4, 8\} \in \Z_{12},\ \ \ |a | = 3.\]
    \item \vspace{-0.2cm}\(\ker\varphi = A_4 \  \mid\ \) Homomorfismo trivial (todo a la identidad), ya considerado en el caso anterior. 
  \end{itemize} 
}

(c) Enuncie el Teorema de Lagrange, y diga  si \(A_4\) tiene subgrupos de orden \(2,3,4,6,8\). 

\textcolor{gray}{
  \(\blacksquare\) \textbf{2.11. (\emph{Lagrange})} Si \(|G| < \infty\) e \( S \leq G \), entonces \(|S|\) divide \(|G|\) e \([G:S] = \nicefrac{|G|}{|S|}\). Esto ya es garantía de que no hay subgrupo orden \(8\nmid 12\). También, justificamos en el ítem (a)  que no hay subgrupo de orden \(6\); 
    \[S_2 := \langle (12)(34)\rangle, \ S_3:= \langle (123)\rangle \text{ e } S_4:= \textbf{V}_4 \leq A_4,\] 
  son subgrupos de los ordenes restantes en la lista. 
}

2. (a) Defina grupo cíclico un diga cuales son los grupos cíclicos a menos de isomorfismo. 

\textcolor{gray}{
  \(G\) és \emph{cíclico} se \(\exists g \in G : G = \langle g \rangle\), existe elemento que genera todo el grupo. Si \(|G| = n < \infty\), entonces 
  \[\varphi : G \cong \Z_n, \text{ donde }\varphi : g \mapsto 1, \]  
  sí el orden fuera infinito tambíén \(\varphi: G \cong \Z\). 
}

(b) Cual es el orden del mayor subgrupo cíclico de \(S_5\)?

\textcolor{gray}{
  La pregunta equivale a cual es el elemento de mayor orden en \(S_5\), si \(\alpha = \beta_1 \ldots \beta_t\) (\(r_j\)-ciclos disjuntos), entonces \(|\alpha| = \operatorname{mcm}\{|\beta_1|, \cdots, |\beta_t|\}\), veamos posibles casos:   
  \[(5); \ (4)(1);\ (3)(2); \ (2)(2)(1); \ (2)(1)(1)(1); \]
  Ahí renta \(\operatorname{mcm}\{3,2\} = 6\), que es mayor que todos los demás. 
}

3. (a) Muestre que los \(3\)-ciclos son conjugados como elementos de \(S_n\) y deduzca \(\exists ! H \leq S_n : [S_n : H] = 2\).  

\textcolor{gray}{
  Sean \((abc)\) e \((xyz) : a < b < c\) e \(x < y < z\), considere: 
  \[\sigma : a \mapsto x; \hspace{0.3cm} \sigma : b \mapsto y; \hspace{0.3cm}   
    \sigma : c \mapsto z;  
  \] 
  e \(\sigma\) fijando las demás entradas, así tenemos \(\sigma (abc) \sigma^{-1} = (\sigma(a) \sigma(b) \sigma(c)) = (xyz)\). Sea \(H \leq A_n : [S_n : H] = 2\), entonces \(H\lhd S_n\) e contiene todos los cuadrados de \(S_n\), en particular todos los \(3\)-ciclos. Por otro lado, \(H \cap A_n \lhd A_n \ \leftrightharpoons \ H \cap A_n = \id \) o \(A_n\), como los \(3\)-ciclos están en \(A_n\), \(H\cap A_n = A_n\), sigue que \(H \subset A_n\) más \(|H| = |A_n|\) pues ambos tienen índice \(2\), finalmente \(H = A_n\).    
}

